% Generated draft: requires mathematical and source review.
\documentclass[11pt]{article}
\usepackage[margin=25mm]{geometry}
\usepackage{fontspec}
\setmainfont{Latin Modern Roman}
\usepackage{amsmath,amssymb,mathtools,unicode-math}
\setmathfont{Latin Modern Math}
\usepackage{enumitem,needspace,xurl,hyperref}
\setlength{\parindent}{0pt}
\setlength{\parskip}{5pt}
\setlength{\emergencystretch}{3em}
\newcommand{\N}{\mathbb{N}}
\newcommand{\Z}{\mathbb{Z}}
\newcommand{\Q}{\mathbb{Q}}
\newcommand{\R}{\mathbb{R}}
\newcommand{\C}{\mathbb{C}}
\newcommand{\F}{\mathbb{F}}
\newcommand{\A}{\mathbb{A}}
\newcommand{\PP}{\mathbb{P}}
\newcommand{\E}{\mathbb{E}}
\newcommand{\eps}{\varepsilon}
\newcommand{\dd}{\,\mathrm{d}}
\begin{document}
\section*{UnsolvedMath ID: 700494}
\textbf{Source problem code:} catalogue entry
\subsection{Definitions and mathematical statement}
For a prime number \(p\) and an integer \(r\) satisfying \(r^2+1\equiv0\pmod p\), define the group
\[
G_{p,r}=\left\langle a,b\ \middle|\ a^2=b^p=(ab)^3=(b^r a b^{-2r}a)^2=1\right\rangle.
\]
Here the presentation denotes the quotient of the free group on \(a,b\) by the normal closure of the displayed relators. Since \(b^p=1\), the presentation depends only on the residue class of \(r\) modulo \(p\).

J. Mennicke asks whether this group is infinite for all but finitely many primes \(p\). The prime domain requires clarification: a permissible \(r\) exists precisely when \(p=2\) or \(p\equiv1\pmod4\). The source also does not explicitly quantify over the possible roots \(r\).

If the intended prime domain is the primes admitting such a root, and the assertion is intended to hold for every root, a precise formulation is
\[
\exists\text{ a finite set }E\text{ of primes}\quad
\forall p\text{ prime with }p\notin E,\quad
\forall r\in\Z,\qquad
r^2+1\equiv0\pmod p\ \Longrightarrow\ |G_{p,r}|=\infty.
\]
This is an explicitly qualified reading; the supplied record does not settle whether every root or a selected root is intended.

The source reports that computer-assisted work proved \(G_p\) finite for \(p\leq17\), subject to the same unresolved admissibility and root conventions. It also states that a solution would have interesting topological applications, without specifying those applications. The supplied registry labels the problem open; neither that label nor the computational report is independently verified here.
\subsection{Short English statement}
Mennicke asks whether the displayed two-generator groups are infinite for all but finitely many eligible primes. Eligibility requires a square root of \(-1\) modulo \(p\); the source leaves the quantification over that root implicit. It reports computer-assisted finiteness for \(p\leq17\) and potential topological applications.
\subsection{Sources}
\begin{itemize}
\item Local UnsolvedMath registry, permanent ID \(700494\), exact source code \texttt{catalogue entry}, title \emph{Kourovka Notebook Problem 8.44}. The supplied complete source record has upstream revision \texttt{d371ee917ac9647fbb06cd6fc529deec6f3dda82} and records status \texttt{open}; its update timestamp is \texttt{2026-10-16T14:00:00Z}.
\item J. Mennicke, Problem 8.44, \emph{The Kourovka Notebook (Unsolved Problems in Group Theory)}, 8th issue (1982), as attributed in the supplied record to the 21st-issue edition (2026), \texttt{arXiv:1401.0300v48}. The supplied collection description names E. I. Khukhro and V. D. Mazurov as editors.
\item The supplied statement reports computer-assisted finiteness for \(p\leq17\) and mentions topological applications. It supplies no separate citation, author, computational details, or application reference for these assertions. The cited Notebook source has not been independently consulted here.
\end{itemize}
\end{document}
